%% ma: L11-B29-0001
%% lop: 11
%% chuong: 8
%% bai: 29
%% dang: 11.29.1
%% loai: TLN
%% muc: TH
%% dap_an: "0,55"
%% nguon: "(NBV)-ĐỀ SỐ 13-MỖI NGÀY 1 ĐỀ THI 2025.pdf (D13, MỖI NGÀY 1 ĐỀ - Đề số 13), Phần III Câu 1"
%% kien_thuc: "Bài 29 (công thức cộng xác suất, biến cố đối); đề dùng thêm khái niệm biến ngẫu nhiên rời rạc và bảng phân bố xác suất không có trong SGK KNTT"
%% trang_thai: cach-ly
%% ly_do: "Đề dùng khái niệm biến ngẫu nhiên rời rạc và bảng phân bố xác suất, không có trong SGK Kết nối tri thức (lời giải chỉ cần công thức cộng và biến cố đối của Bài 29). Gợi ý sửa đề: thay bằng bảng xác suất của các biến cố E_0, E_1, ... (xảy ra đúng k cuộc gọi), không nhắc biến ngẫu nhiên."
%% ghi_chu: "Đáp án gốc 0,55 đúng (1 - 0,25 - 0,2 = 0,55). Đã bỏ ảnh minh họa điện thoại."
\begin{ex}
Số cuộc điện thoại gọi đến một trung tâm cứu hộ trong khoảng thời gian từ $12$ giờ đến $13$ giờ là một biến ngẫu nhiên rời rạc $X$ có bảng phân bố xác suất như sau:
\begin{center}
\begin{tabular}{|c|c|c|c|c|c|c|}\hline
$X$ & $0$ & $1$ & $2$ & $3$ & $4$ & $5$\\\hline
$P$ & $0{,}25$ & $0{,}2$ & $0{,}15$ & $0{,}15$ & $0{,}13$ & $0{,}12$\\\hline
\end{tabular}
\end{center}
Tính xác suất để xảy ra ít nhất $2$ cuộc gọi đến trung tâm cứu hộ đó.
\shortans{0,55}
\loigiai{Gọi $E$: ``xảy ra ít nhất $2$ cuộc gọi''. Biến cố đối $\overline E$: ``xảy ra nhiều nhất $1$ cuộc gọi'' là hợp của hai biến cố xung khắc $\{X=0\}$ và $\{X=1\}$, nên $P(\overline E)=0{,}25+0{,}2=0{,}45$. Vậy $P(E)=1-0{,}45=0{,}55$.}
\end{ex}
